% PLAN:N8
% NODES:
\section{Conclusion}\label{sec:N8}

% N8-p1
Spectral Surface Generation is fixed by its input-output contract. A quantum device samples the spectral signal $G(t)$ at the Gauss-Hermite grid times $t_k$ in one acquisition campaign, banking $\hat{G}(t_k)$. A classical layer reconstructs the Gaussian-filtered spectral surface $\hat{Z}_\tau(s)$, the generating surface and the response channels - filtered mean and variance, higher cumulants, within-window ratios - at post-acquisition $(s,\tau)$. Hats mark finite-record estimators - \emph{exact} is reserved for the continuous transform and the algebraic identities. The Fourier reconstruction identity makes each read and each coefficient insertion a coefficient update on the stored record - no further quantum evolution. The $\tau$ direction obeys a forward-stable heat-kernel scale space, and the classical reconstruction is costed under a named query model.

% N8-p2
One acquisition campaign supports classical reconstruction of $\hat{Z}_\tau(s)$ and its response channels only over the finite-record envelope - the $(s,\tau)$ domain bounded by time reach from the Gaussian aperture, node retention from the Gauss-Hermite gate and the shot-noise floor. Shot noise propagates through the linear reconstruction with a derived covariance for shared-record ratio estimators, whose gain is a constant factor. The standing resource limit is the visibility wall: an absolute read of exponentially small spectral weight costs exponentially many shots - a cost bound, not an estimator defect - so every readout claim is conditioned on its window clearing the shot floor.

% N8-p3
On the controlled subset-sum family every reconstructed quantity has an independently computable validation reference: at $V = 0$ the plus-state overlap factorises exactly and the degeneracy profile $g(T)$ is an exact Fourier coefficient of the record. The reconstruction agreed with that reference. The multiplicity fingerprint responded at order $J^2$ within its stated forward model - it reads the compressed perturbation's within-cluster variance, not $g$ itself - with the $g = 1$ null channel at its finite empirical floor. The integer decoder reported coverage and conditional accuracy under its erasure rule inside a nonempty measured operating window. Every numerical statement names the method behind it - exact diagonalisation at $n \le 12$, synthetic finite-shot sampling at stated budgets $K$, TEBD at $n \le 20$. No quantum hardware data exists. Tested sizes are proof of principle and the reconstruction is validated under trusted assumptions rather than certified. The benchmark inherits no worst-case hardness - instances at the production weight scale are tractable by dynamic programming.

% N8-p4
What the benchmark does not settle is hardness. For the tested distinct-weight instances inside the readable window the dynamics remained classically compressible - a finite-size statement, not a hardness no-go. On the planted equal-weight family readability coexisted with bond-dimension growth consistent with exponential over the tested range - a guide, not an established law. The hardness of the planted family is open. The next technical question is its formalisation: whether the planted construction defines an instance family whose classical simulation cost grows as the finite-size guide suggests and can be tied to a named hardness assumption, while the banked record keeps supplying readable diagnostics. No advantage claim rests on the answer.
